Задача оценки истинных параметров математической модели встречается во многих научных и прикладных областях: 
    генетике [fast-eps-source], физике [DINGO-source], финансах и др. Она становится особенно сложной, когда 
    невозможно аналитически описать распределение наблюдений при заданных параметрах. В этой работе мы, 
    следуя статье [DINGO-source], предлагаем модификацию метода Flow Matching Posterior Estimation
    и сравниваем его с другими методами на простых задачах, а в конце применем его для оценки параметеров 
    модели стохастической волатильности.

    % Итак, наша обычно задача состоит из следующих частей
    % \begin{enumerate}
    %     \item Параметрическая модель некоторого процесса
    %     \item Симулятор 
    %     \item Данные наблюдений/симуляций
    % \end{enumerate} 
    % Цель обычно состоит в том, чтобы по наблюдениям подобрать наиболее согласованные с данными параметры модели

    % В настоящее время классичесие методы используют нейросети для аппроксимации искомого распределения параметров.
    % На выходе мы получаем функцию, которая по наблюдениям выдает параметры.
    % Но зачастую нам требуется знать параметры только для одного наблюдения, а не для всех возможных.

    % % Эта задача называется Sequential Neural Posterior Estimation

    % В этой работе мы сосредоточимся на Sequential модификации метода FMPE [DINGO-source], 
    % который является state-of-the-art на данный момент.

    С математической точки зрения наша задача формулируется следующим образом. 
    Будем обозначать $\theta \in \mathbb{R}^m$ --- вектор параметров модели, а $x \in \mathbb{R}^n$ --- наблюдаемые/симулированные данные
    или, чаще, их малоразмерная аппроксимация, выжимка. Из-за стохастической природы моделей удобно считать как $x$, 
    так и $\theta$ случайными величинами. 
    
    Далее, нам известно, возможно неявно, \textit{правдоподобие} $p(x \mid \theta)$ - распределение наблюдений $x$, которые мы можем получить при условии, 
    что модель инициализирована параметрами $\theta$. Также мы имеем вольность задать \textit{априорное} распределение $p(\theta)$ 
    - наше общее представление о том, в каких пределах могут изменяться параметры модели. 

    Имея все это, мы хотим узнать \textit{апостериорное} распределение $p(\theta \mid x)$ --- распределенине параметров, 
    которые могли бы дать нам такое же наблюдение $x$. Тогда по формуле байеса имеем 
    $$
    p(\theta \mid x) = \frac{p(x \mid \theta) p(\theta)}{p(x)}
    $$
    
    \paragraph{\textbf{Markov Chain Monte Carlo}} Таким образом, коль скоро нам известны $p(x \mid \theta), p(\theta)$, можно считать, 
    что $p(\theta \mid x) \propto p(x \mid \theta) p(\theta)$. Если бы мы имели аналитическое выражение для последнего, 
    то могли бы сэмплировать из апостериорного распределения методами Markov Chain Monte Carlo (MCMC, \cite{MCMC}). 
    Их суть заключается в том, чтобы построить марковскую цепь, стационарное распределение которой будет совпадать 
    с заданным. Тогда, пройдя по ней достаточное число шагов, мы получим выборку из нужного распределения.
    
    Помимо того, что нам нужно знать явное выражение для ненормированной плотности, 
    у этого метода есть другие минусы, такие как: вычислительная затратность, сложность диагностики сходимости,
    чувствительность к настройкам и проблемы в больших размерностях.
    
    \paragraph{\textbf{Approximate Bayesian Computation}} В случае когда мы не можем аналитически выразить правдоподобие, 
    приходится действовать другими методами. В литературе их объединяют под общим названием ABC (Approximate Bayesian Computation \cite{ABC}) методов. 
    Классические алгоритмы ABC основаны на предоположении о том, что если наблюдения похожи, то и соответствующие параметры тоже.
    Более формально, они аппроксимируют истинное апостериорное $p(\theta \mid x_o)$ с помощью $p(\theta \mid \lvert x_o - x \rvert < \varepsilon)$
    
    Самый простой метод, Rejection ABC, работает так:
    \begin{enumerate}
        \item Дано: наблюдение $x_o$, хотим получить выборку из $p(\theta \mid x_o)$
        \item Генерируем $\theta^* \sim p(\theta)$
        \item Симулируем $x \sim p(x \mid \theta^*)$
        \item Если $\lvert x_o - x \rvert < \varepsilon$, то добавляем $\theta^*$ в выборку, иначе повторить 2.
    \end{enumerate} 
    Из очевидных минусов --- это проклятие размерности и неэффективность вычислений. 
    Для вычисления апостериора только в одной точке $x_o$ нам требуется провести тысячи симуляций, 
    большая часть из которых будет отметена, причем с увеличением размености это число растет экспоненциально.

    \paragraph{\textbf{Neural Posterior Estimation}} В последние годы, с развитием машинного обучения, 
    пришли и новые методы оценки апостериора. В отличие от классических алгоритмов, где апостериор вычислялся 
    только эмпирически, методы NPE (Neural Posterior Estimation, \cite{NPE}) позволяют аппроксимировать функцию
    плотности с помощью нейросети. Часто в качестве универсального аппроксиматора используют \textit{нормализующие потоки} (Normalizing Flows, \cite{NormFlows}),
    хорошо показавшие себя в задачах физики гравитационных волн \cite{GravWaves}. Нормализующие потоки отображают 
    сэмплы из одного (обычно, гауссовского) распределения в другое распределение серией последовательных несложных преобразований, 
    параметризованых нейросетью. 
    
    Минус такого подхода в том, что эти преобразования должны быть обратимы 
    и иметь определенный вид чтобы сеть можно было эффективно обучать. К тому же, такая архитектура плохо масштабируема 
    для нужд больших и сложных экспериментов. 
    
    Однако, самый большой и главный плюс в сравнении с классическими методами --- это то, что теперь, у нас на выходе
    не просто выборка из нужного распределения, а функция, которая каждому наблюдению $x$ сопоставляет распределение $p(\theta \mid x)$.
    Таким образом, один раз обучив такую модель, мы получим $p(\theta \mid x)$ для всех возможных $x$, а не только одного 
    --- в литературе это называется амортизированным инференсом (Amortized Inference, \cite{AmortInfr}). 
    
    В контексте отображения одного распределения в другое эта задача родственна задаче генеративного ИИ. 
    А в последние годы, в этой области доминируют диффузионные модели \cite{DiffModels}. Принцип их работы 
    основан на том же преобразовании одного распределения в другое, только теперь сэмплы двигаются по траекториям, 
    которые являются решениями СДУ, заданного векторным полем. И нейросетью здесь параметризовано это самое поле.
    
    Детерминированным аналогом диффузионных моделей являются Flow Matching модели \cite{FlowMatching}. В данном случае, 
    вместо стохастического дифференциального уравнения используется обыкновенное. На них мы и сосредоточимся в этой работе.
    

    
    % \paragraph{Normalizing Flows} Нормализующий поток (Normalizing Flow, \cite{NormFlow}) задает условное распределение параметров $q(\theta \mid x), \theta \in \mathbb{R}^n, \theta \in \mathbb{R}^m$, 
    % которое определяется как 
    % \[
    % q(\theta \mid x) = [\psi_x]_# q_0(\theta) = q_0(\psi_x^{-1}(\theta)) \det \Big| \partial_\theta \psi_x^{-1}(\theta) \Big|,
    % \]
    % где $\psi_x: \mathbf{R}^n \to \mathbb{R}^n$

    
    % \paragraph{Flow Matching}
    